\title{Direct Preference Optimization: Your Language Model is Secretly a Reward Model}

\begin{abstract}
Although large-scale unsupervised language models (LMs) acquire extensive world knowledge and some reasoning capabilities, precisely steering their behavior remains challenging due to the fully unsupervised nature of their training. Current approaches to achieve such control involve gathering human annotations on the relative quality of model outputs and fine-tuning the unsupervised LM to align with these preferences, often through reinforcement learning from human feedback (RLHF). Nevertheless, RLHF is a complicated and frequently unstable process that first trains a reward model to represent human preferences, then fine-tunes the large unsupervised LM via reinforcement learning to maximize this predicted reward while avoiding excessive deviation from the original model. In this work, we propose a novel parameterization of the reward model in RLHF that allows direct extraction of the corresponding optimal policy in closed form, enabling us to address the standard RLHF problem using only a straightforward classification loss. The resulting method, which we term \textit{Direct Preference Optimization} (DPO), is stable, effective, and computationally efficient, removing the necessity for sampling from the LM during fine-tuning or extensive hyperparameter tuning. Our experiments demonstrate that DPO can fine-tune LMs to align with human preferences as well or better than existing techniques. Importantly, fine-tuning with DPO surpasses PPO-based RLHF in controlling generation sentiment, and equals or improves response quality in summarization and single-turn dialogue, all while being considerably simpler to implement and train.
\end{abstract}

\section{Introduction}
Large-scale unsupervised language models (LMs) trained on extensive datasets develop remarkable capabilities~\citep{chowdhery2022palm, brown2020language, touvron2023llama,bubeck2023sparks}. Nonetheless, these models learn from data generated by humans who have diverse objectives, priorities, and skill levels. Some of these objectives and skills may not be desirable to replicate; for instance, while we want our AI coding assistant to \textit{recognize} common programming errors to correct them, when generating code, we prefer to guide the model toward the (potentially infrequent) high-quality coding skills embedded in its training data. Similarly, we might want our language model to \textit{be aware} of a widespread misconception held by 50\% of people, yet we certainly do not want the model to assert this misconception as true in half of the queries about it. Put differently, selecting the model’s \emph{preferred responses and behaviors} from its extensive \textit{knowledge and capabilities} is vital for developing AI systems that are safe, effective, and controllable \citep{ouyang2022training}. Although existing approaches generally steer LMs to align with human preferences using reinforcement learning (RL), we will demonstrate that the RL-based objective utilized in these methods can be exactly optimized via a straightforward binary cross-entropy objective, significantly streamlining the preference learning process.

Broadly speaking, current methods embed desired behaviors into a language model through carefully curated datasets of human preferences that reflect the types of behavior humans consider safe and useful. This preference learning phase follows an initial stage of large-scale unsupervised pre-training on an extensive text corpus. The simplest approach to preference learning involves supervised fine-tuning on human demonstrations of high-quality responses; however, the most effective techniques are reinforcement learning from human (or AI) feedback (RLHF/RLAIF; \citep{christiano2017deep,bai2022constitutional}). RLHF methods train a reward model on a dataset of human preferences and then employ RL to optimize the language model policy to generate responses with high reward while avoiding excessive deviation from the original model. Although RLHF yields models with strong conversational and coding skills, the RLHF pipeline is substantially more complex than supervised learning, requiring training multiple language models and performing in-loop sampling from the LM policy, which results in considerable computational overhead.

In this paper, we demonstrate a method to directly optimize a language model according to human preferences, without relying on explicit reward modeling or reinforcement learning. We introduce \textit{Direct Preference Optimization} (DPO), an algorithm that implicitly targets the same objective as current RLHF algorithms (reward maximization constrained by KL-divergence) but is easier to implement and simpler to train. Conceptually, the DPO update boosts the relative log probability of preferred responses compared to dispreferred ones, while applying a dynamic, example-specific importance weight that avoids the model degradation observed with a naive probability ratio objective. Similar to existing approaches, DPO depends on a theoretical preference model (e.g., the Bradley-Terry model; \cite{bradley1952rankanalysis}) to assess how well a reward function corresponds to empirical preference data. However, whereas prior methods employ the preference model to define a loss for training a reward model, followed by policy optimization on the learned reward, DPO applies a change of variables to express the preference loss directly as a function of the policy. Consequently, given a dataset of human preferences over model outputs, DPO can optimize the policy via a straightforward binary cross entropy objective, yielding the optimal policy for an implicit reward function inferred from the preference data.

Our primary contribution is Direct Preference Optimization (DPO), a straightforward algorithm that trains language models from preferences without reinforcement learning. Our experiments demonstrate that DPO performs at least as well as established techniques, including PPO-based RLHF, when learning from preferences across tasks such as sentiment adjustment, summarization, and dialogue, using language models with up to 6 billion parameters.

\section{Related Work}

Self-supervised language models with increasing size are capable of completing certain tasks either zero-shot \citep{radford2019language} or through few-shot prompting \citep{gpt3,megatron,chowdhery2022palm}. Nevertheless, their effectiveness on downstream tasks and alignment with user intentions can be substantially enhanced by fine-tuning on datasets consisting of instructions and human-generated completions \citep{mishra-etal-2022-cross,sanh2022multitask,chung2022scaling,thoppilan2022lamda}. This process of `instruction-tuning' allows LLMs to generalize to instructions beyond those seen during tuning and generally improve their usability \citep{chung2022scaling}. Although instruction tuning has been successful, \textit{relative} human evaluations of response quality are often simpler to obtain than expert demonstrations. Consequently, later studies have fine-tuned LLMs using datasets of human preferences, enhancing performance in areas such as translation \citep{kreutzer-etal-2018-reliability}, summarization \citep{stiennon2022learning,ziegler2020finetuning}, storytelling \citep{ziegler2020finetuning}, and instruction-following \citep{ouyang2022training,ramamurthy2023is}. These approaches initially train a neural network reward model to align with preference datasets under a preference framework like the Bradley-Terry model \citep{bradley1952rankanalysis}, then fine-tune the language model to maximize this reward via reinforcement learning techniques, typically REINFORCE \citep{williams1992reinforce}, proximal policy optimization (PPO; \cite{schulman2017proximal}), or related variants \citep{ramamurthy2023is}. A related line of research employs LLMs fine-tuned with human feedback on instructions to generate synthetic preference data focusing on specific attributes such as safety or harmlessness \citep{bai2022constitutional}, relying only on weak human supervision in the form of textual rubrics for the LLM's annotations. These methods reflect the convergence of two research areas: one focused on training language models with reinforcement learning for diverse objectives~\citep{Ranzato2015SequenceLT,paulus2018a,wu2018learning} and another centered on general frameworks for learning from human preferences \citep{christiano2017deep,kupcsik2018learning}. Despite the advantages of using relative human preferences, fine-tuning large language models with reinforcement learning remains a significant practical hurdle; this paper proposes a theoretically grounded method for optimizing relative preferences without relying on RL.

Outside the realm of language, policy learning from preferences has been explored within both bandit and reinforcement learning frameworks, with multiple methods proposed. Contextual bandit learning that relies on preferences or rankings of actions instead of explicit rewards is termed a contextual dueling bandit (CDB; \cite{yue2012karmed,dudik2015contextual}). When absolute rewards are unavailable, theoretical analyses of CDBs replace the concept of an optimal policy with a \textit{von Neumann winner}, defined as a policy that has an expected winning probability of at least 50\% against \textit{any} other policy \citep{dudik2015contextual}. Nonetheless, in CDB frameworks, preference labels are provided in an online manner, whereas learning from human preferences generally involves training from a static batch of offline preference-labeled action pairs \citep{yan2022human}. Similarly, \textit{preference-based RL} (PbRL) utilizes binary preferences generated by an \textit{unknown} ‘scoring’ function instead of direct reward signals \citep{BusaFekete2014,ruiz2023dueling}. Various PbRL algorithms have been developed, including those capable of leveraging off-policy preference data, but they typically involve first explicitly inferring the latent scoring function (i.e., the reward model) followed by its optimization \citep{jain2013learning,BusaFekete2014,christiano2017deep,sadigh2017active,kupcsik2018learning}. In contrast, we propose a unified policy learning approach that directly optimizes a policy to satisfy the given preferences.

\section{Preliminaries}\label{section:prelims}

We provide a review of the RLHF pipeline as described by \citeauthor{ziegler2020finetuning} and subsequent works \citep{stiennon2022learning, bai2022training, ouyang2022training}. This pipeline generally involves three stages: 1) supervised fine-tuning (SFT); 2) preference sampling and reward modeling; and 3) reinforcement learning optimization.

\textbf{SFT}: The RLHF process typically starts by fine-tuning a pre-trained language model on high-quality labeled data for the target downstream tasks (e.g., dialogue, summarization) using supervised learning, resulting in a model denoted as $\pi^\text{SFT}$.

\textbf{Reward Modeling Phase}: In the subsequent stage, the SFT model is prompted with inputs $x$ to generate pairs of responses $(y_1, y_2) \sim \pi^\text{SFT}(y \mid x)$. These response pairs are then evaluated by human annotators who express preferences between them, represented as $y_w \succ y_l \mid x$, where $y_w$ and $y_l$ refer to the preferred and less preferred outputs among $(y_1, y_2)$, respectively. The preferences are assumed to be generated by some underlying reward function $r^*(y, x)$, which is unobserved. Multiple methods exist to model such preferences, with the Bradley-Terry (BT) \cite{bradley1952rankanalysis} model being a widely used choice (although more general models like the Plackett-Luce ranking models \citep{plackett1975analysis, luce2012individual} can also be applied if multiple ranked answers are available). According to the BT model, the human preference distribution $p^*$ can be expressed as:

\begin{equation}\label{eq:bradley-terry}
    p^*(y_1\succ y_2 \mid x)=\frac{\exp\left(r^*(x, y_1)\right)}{\exp\left(r^*(x, y_1)\right) + \exp\left(r^*(x, y_2)\right)}.
\end{equation}

Given access to a fixed dataset of comparisons $\mathcal{D}=\bigl\{x^{(i)}, y_w^{(i)}, y_l^{(i)}\bigr\}_{i=1}^N$ drawn from $p^*$, we can model a reward function $r_{\phi}(x, y)$ with parameters estimated by maximizing the likelihood. By framing the task as binary classification, the negative log-likelihood loss can be expressed as:

\begin{equation}\label{eq:reward_model}
    \mathcal{L}_R(r_{\phi}, \mathcal{D}) = -\mathbb{E}_{(x, y_w, y_l)\sim \mathcal{D}}\bigl[\log \sigma(r_{\phi}(x, y_w)- r_{\phi}(x, y_l))\bigr]
\end{equation}

where $\sigma$ denotes the logistic function. In language models, the function $r_{\phi}(x, y)$ is commonly initialized from the SFT model $\pi^\text{SFT}(y \mid x)$ by adding a linear layer on top of the final transformer layer that outputs a scalar reward prediction \cite{ziegler2020finetuning}. To reduce reward variance, previous works normalize rewards to ensure $\mathbb{E}_{x,y\sim \mathcal{D}}\left[r_\phi(x, y)\right] = 0$ for all $x$.

\textbf{RL Fine-Tuning Phase}: In the reinforcement learning phase, the learned reward function guides the language model's updates. Following earlier studies~\citep{jaques2017sequence, jaques2020human}, this optimization is formulated as

\begin{equation}\label{eq:RL}
\max_{\pi_{\theta}}  \mathbb{E}_{x\sim \mathcal{D}, y\sim \pi_{\theta}(y \mid x)}\bigl[r_{\phi}(x, y)\bigr] - \beta\mathbb{D}_{\textrm{KL}}\bigl[\pi_{\theta}(y\mid x)\mid \mid \pi_\text{ref}(y\mid x)\bigr],
\end{equation}

where $\beta$ regulates the extent of deviation from the reference policy $\pi_\text{ref}$, which is the initial SFT model $\pi^\text{SFT}$. 
In practical scenarios, the policy $\pi_\theta$ is initialized to $\pi^\text{SFT}$. This regularization is crucial to prevent the policy from straying too far from the distribution where the reward model is reliable, and to maintain diversity in generation while avoiding collapse to a few high-reward outputs. Since language generation is discrete, this objective is non-differentiable and usually optimized via reinforcement learning. The typical method \citep{ziegler2020finetuning, stiennon2022learning, bai2022training, ouyang2022training} involves defining the reward as ${r(x, y) = r_{\phi}(x, y) -\beta (\log \pi_{\theta}(y\mid x) - \log \pi_\text{ref}(y\mid x))}$ and maximizing it using PPO \cite{schulman2017proximal}. 

\section{Direct Preference Optimization}\label{sec:DPO}

To address the difficulties of applying reinforcement learning methods to large-scale tasks like fine-tuning language models, we propose a straightforward approach to optimize policies directly using preference data. In contrast to prior RLHF techniques that first learn a reward function and then optimize it via RL, our method employs a specific reward model parameterization that allows the optimal policy to be derived analytically, bypassing the need for an RL training loop. Our central insight is to use an exact mapping from reward functions to their optimal policies, transforming a loss defined over reward functions into one defined over policies. This change of variables eliminates the necessity of fitting an explicit reward model while still optimizing under standard human preference frameworks such as the Bradley-Terry model. Essentially, the policy network serves as both the language model and the (implicit) reward function.

\textbf{Deriving the DPO objective.} We begin with the same reinforcement learning (RL) objective as previous studies, Eq.~\ref{eq:RL}, involving a general reward function $r$. Following earlier work~\citep{peters2007reinforcement, peng2019advantage, korbak2022reinforcement, go2023aligning}, it is straightforward to demonstrate that the optimal solution to the KL-constrained reward maximization problem in Eq.~\ref{eq:RL} has the form:

\begin{equation}\label{eq:op_policy}
    \pi_r(y\mid x) = \frac{1}{Z(x)}\pi_\text{ref}(y\mid x)\exp\left(\frac{1}{\beta}r(x, y)\right),
\end{equation}

where the partition function $Z(x) = \sum_{y}\pi_\text{ref}(y\mid x)\exp\left(\frac{1}{\beta}r(x, y)\right)$. A full derivation can be found in Appendix \ref{app:derivation1}. Even when employing the maximum likelihood estimate $r_{\phi}$ of the true reward function $r^*$, estimating the partition function $Z(x)$ remains computationally expensive~\citep{korbak2022reinforcement, go2023aligning}, making this form difficult to use practically. Nonetheless, we can rearrange Eq.~\ref{eq:op_policy} to rewrite the reward function in terms of the corresponding optimal policy $\pi_r$, the reference policy $\pi_\text{ref}$, and the unknown partition function $Z(\cdot)$. Specifically, by taking the logarithm of both sides of Eq.~\ref{eq:op_policy} and performing algebraic manipulation, we arrive at:

\begin{equation}\label{eq:main_eq}
    r(x,y) = \beta \log \frac{\pi_r(y\mid x)}{\pi_\text{ref}(y\mid x)} + \beta \log Z(x).
\end{equation}

This reparameterization can be applied to the ground-truth reward $r^*$ and its corresponding optimal policy $\pi^*$. Crucially, the Bradley-Terry model depends solely on the reward differences between two completions, i.e., ${p^*(y_1 \succ y_2 \mid x) = \sigma(r^*(x, y_1) - r^*(x, y_2))}$. Substituting the expression for $r^*(x,y)$ from Eq.~\ref{eq:main_eq} into the preference model of Eq.~\ref{eq:bradley-terry}, the partition function terms cancel out, allowing us to express the human preference probability entirely through the optimal policy $\pi^*$ and the reference policy $\pi_\text{ref}$. Therefore, the optimal RLHF policy $\pi^*$ under the Bradley-Terry model adheres to the preference model:

\begin{equation}\label{eq:objective}
    p^*(y_1\succ y_2 \mid x) = \frac{1}{1 + \exp\left(\beta \log \frac{\pi^*(y_2\mid x)}{\pi_\text{ref}(y_2\mid x)} - \beta \log \frac{\pi^*(y_1\mid x)}{\pi_\text{ref}(y_1\mid x)}\right)}.
\end{equation}

The full derivation is provided in Appendix~\ref{app:derivation2}. While Eq.~\ref{eq:objective} is based on the Bradley-Terry model, analogous formulations can be derived for more general Plackett-Luce models~\citep{plackett1975analysis, luce2012individual}, as detailed in Appendix~\ref{app:plackett_luce_models}.

Having expressed the probability of human preference data in terms of the optimal policy rather than the reward model, we can now define a maximum likelihood objective for a parameterized policy $\pi_\theta$. Similar to the reward modeling framework (i.e., Eq.~\ref{eq:reward_model}), the policy objective is:

\begin{equation}\label{eq:optimum_model}
    \mathcal{L}_\text{DPO}(\pi_{\theta}; \pi_\text{ref}) = -\mathbb{E}_{(x, y_w, y_l) \sim \mathcal{D}}\left[\log \sigma \left(\beta \log \frac{\pi_{\theta}(y_w\mid x)}{\pi_\text{ref}(y_w\mid x)} - \

In this manner, we fit an implicit reward through a different parameterization, where the optimal policy corresponds directly to $\pi_\theta$. Additionally, because our method is equivalent to fitting a reparameterized Bradley-Terry model, it benefits from certain theoretical guarantees, such as consistency given appropriate assumptions about the distribution of preference data \cite{bong2022generalized}. Section~\ref{sec:theory} provides further discussion on the theoretical aspects of DPO in comparison to related work.

\textbf{What effect does the DPO update have?} To gain a mechanistic insight into DPO, it is helpful to examine the gradient of the loss function $\mathcal{L}_\text{DPO}$. The gradient with respect to parameters $\theta$ can be expressed as:

\begin{multline*}\label{eq:gradient}
    \nabla_\theta \mathcal{L}_\text{DPO}(\pi_\theta;\pi_\text{ref}) = \\ -\beta\mathbb{E}_{(x, y_w, y_l) \sim \mathcal{D}} \bigg[\underbrace{\sigma(\hat{r}_\theta(x, y_l) - \hat{r}_\theta (x, y_w))}_\text{greater weight when reward prediction is incorrect}\bigg[\underbrace{\nabla_\theta\log \pi(y_w \mid x)}_\text{increase likelihood of $y_w$} - \underbrace{\nabla_\theta\log\pi(y_l \mid x)}_\text{decrease likelihood of $y_l$}\bigg]\bigg],
\end{multline*}

where $\hat{r}_\theta(x, y) = \beta \log \frac{\pi_\theta(y \mid x)}{\pi_\text{ref}(y \mid x)}$ denotes the reward implicitly defined by the language model $\pi_\theta$ relative to the reference model $\pi_\text{ref}$ (further details in Section~\ref{sec:theory}). Intuitively, the gradient of $\mathcal{L}_\text{DPO}$ adjusts the probabilities by increasing the likelihood of preferred completions $y_w$ while reducing that of less favored completions $y_l$. Crucially, the samples are weighted according to how much higher the implicit reward $\hat{r}_\theta$ rates the disfavored completion, modulated by $\beta$. In other words, the weighting reflects the extent to which the implicit reward model incorrectly ranks the outputs, factoring in the strength of the KL regularization. Our experiments highlight the importance of this weighting, as a straightforward variant lacking this weighting term can lead to model collapse (see Appendix Table~\ref{tab:unlikelihood_generations}).

\textbf{DPO overview.} The typical DPO process consists of the following steps: 1) For each prompt $x$, sample completions $y_1, y_2 \sim \pi_\text{ref}(\cdot \mid x)$, label them using human preferences to create the offline preference dataset $\mathcal{D} = \{x^{(i)}, y_w^{(i)}, y_l^{(i)}\}_{i=1}^N$, and 2) train the language model $\pi_\theta$ by minimizing the loss $\mathcal{L}_\text{DPO}$ given the reference policy $\pi_\text{ref}$, dataset $\mathcal{D}$, and target parameter $\beta$. 
In practical scenarios, it is preferable to reuse existing publicly available preference datasets instead of generating new samples and collecting human feedback. Because these preference datasets are generated by sampling from $\pi^\text{SFT}$, we set $\pi_\text{ref} = \pi^\text{SFT}$ whenever possible. If $\pi^\text{SFT}$ is not accessible, we initialize $\pi_\text{ref}$ by maximizing the likelihood of the preferred completions ${(x, y_w)}$, i.e., ${\pi_\text{ref} = \argmax_{\pi} \mathbb{E}_{x, y_w \sim \mathcal{D}}[\log \pi(y_w \mid x)]}$. This approach reduces the distributional mismatch between the unavailable true reference distribution and the $\pi_\text{ref}$ used in DPO. Additional details regarding implementation and hyperparameters are provided in Appendix~\ref{app:implementation}.

\section{Theoretical Analysis of DPO} This section offers a deeper interpretation of the DPO method, presents theoretical justification, and connects the benefits of DPO to limitations of actor-critic algorithms employed for RLHF, such as PPO~\cite{schulman2017proximal}.

\label{sec:theory}

\subsection{Your Language Model Is Secretly a Reward Model} DPO manages to circumvent both explicit reward fitting and reinforcement learning by optimizing a single maximum likelihood objective. Notice that the optimization objective in Eq. \ref{eq:main_eq} corresponds to a Bradley-Terry model with reward parameterization $r^*(x, y) = \beta \log\frac{\pi^*_\theta(y \mid x)}{\pi_\text{ref}(y \mid x)}$, and we optimize the parametric model $\pi_{\theta}$ equivalently to reward model training described in Eq. \ref{eq:reward_model}, subject to a change of variables. In this section, we develop the theoretical foundation for this reparameterization, demonstrate that it does not restrict the class of reward models that can be learned, and show that it enables exact recovery of the optimal policy. We start by defining an equivalence relation between reward functions.

\begin{definition}
We define two reward functions $r(x, y)$ and $r'(x, y)$ as equivalent if and only if there exists some function $f$ such that ${r(x, y)-r'(x, y) = f(x)}$.     
\end{definition}

It is straightforward to verify that this relation is indeed an equivalence relation, which divides the set of reward functions into distinct equivalence classes. The following two lemmas can be stated:

\begin{lemma}\label{lemma:same_prefrence} Under the Plackett-Luce framework, and specifically the Bradley-Terry model, two reward functions belonging to the same equivalence class generate the same preference distribution.
\end{lemma}

\begin{lemma}\label{lemma:same_policy}
    Any two reward functions within the same equivalence class lead to the identical optimal policy in the constrained RL problem.
\end{lemma}

The proofs are direct and are provided in Appendix \ref{app:lemma1}. The first lemma highlights a known issue of under-specification in the Plackett-Luce model family \cite{plackett1975analysis}. Due to this under-specification, it is common to enforce additional identifiability constraints to secure guarantees for the MLE estimations from Eq. \ref{eq:reward_model} \cite{bong2022generalized}. The second lemma asserts that all reward functions from a single equivalence class yield the same optimal policy, so our ultimate goal is only to recover some reward function within that optimal class. The following theorem is proven in Appendix~\ref{app:thm1}:

\begin{theorem}\label{thm:main}
    Under reasonable assumptions, every reward class compatible with the Plackett-Luce model (and in particular the Bradley-Terry model) can be expressed by the reparameterization ${r(x, y) = \beta \log \frac{\pi(y\mid x)}{\pi_\text{ref}(y\mid x)}}$ for some model $\pi(y\mid x)$ and a given reference model $\pi_\text{ref}(y \mid x)$.
\end{theorem}

\begin{sproof}
    Take any reward function $r(x, y)$, which induces an associated optimal model $\pi_r(y \mid x)$ defined by Eq. \ref{eq:op_policy}. We will demonstrate that a reward function equivalent to $r$ can be expressed using the stated reparameterization. Define the projection operator $f$ as  
\begin{equation}
    f(r; \pi_\text{ref}, \beta)(x, y) = r(x, y) - \beta\log\sum_{y}\pi_\text{ref}(y\mid x)\exp\left(\frac{1}{\beta}r(x, y)\right)
\end{equation}
This operator $f$ essentially normalizes the reward function by subtracting the logarithm of the partition function related to $\pi_r$. Because this normalization term depends only on the prefix $x$, $f(r; \pi_\text{ref}, \beta)(x, y)$ remains in the equivalence class of $r(x, y)$. Finally, by substituting $r$ with the right-hand side of Eq.~\ref{eq:main_eq} (which is valid for any reward function), we obtain $f(r; \pi_\text{ref}, \beta)(x, y) = \beta \log \frac{\pi_r(y\mid x)}{\pi_\text{ref}(y\mid x)}$. Thus, the projection $f$ yields a member of $r$’s equivalence class with the desired form, showing that the proposed reparameterization does not reduce the generality of our reward model.
\end{sproof}

An alternative perspective on Theorem~\ref{thm:main} is that it precisely identifies which reward function within each equivalence class the DPO reparameterization chooses, specifically, the reward function that satisfies:

\begin{equation}\label{eq:lag_p}
     \sum_{y}\underbrace{\pi_\text{ref}(y\mid x)\exp\left(\frac{1}{\beta}r(x, y)\right)}_{=\pi(y\mid x)\text{, using Thm.~\ref{thm:main} reparam.}} = 1,
\end{equation}

meaning that $\pi(y\mid x)$ constitutes a proper probability distribution (with positive probabilities summing to one). Notably, from Eq.~\ref{eq:op_policy}, we recognize that Eq.~\ref{eq:lag_p} represents the partition function of the optimal policy derived from the reward function $r(x, y)$. The crucial insight behind the DPO algorithm is that by enforcing certain constraints on the otherwise under-determined Plackett-Luce family of preference models (and particularly the Bradley-Terry subclass), we maintain the representable reward model class while ensuring that the optimal policy in Eq.~\ref{eq:op_policy} becomes analytically manageable for every prompt $x$.

\subsection{Instability of Actor-Critic Algorithms} Our framework also enables an analysis of the instabilities found in typical actor-critic algorithms employed in RLHF, such as PPO. Focusing on the RL fine-tuning phase described in Section \ref{section:prelims}, we can relate this to the control as inference framework \cite{levine2018reinforcement} applied to the constrained RL problem stated in \ref{eq:RL}. Consider a parameterized policy model $\pi_{\theta}(y\mid x)$; the objective is to minimize the divergence $\mathbb{D}_{\text{KL}}[\pi_{\theta}(y|x) \mid \mid \pi^*(y\mid x)]$, where $\pi^*$ is the optimal policy from Eq.~\ref{eq:optimum_model} induced by the reward $r_{\phi}(y, x)$. Through algebraic manipulation, this yields the optimization objective:

\begin{equation}\label{eq:AC}
    \max_{\pi_{\theta}}\mathbb{E}_{\pi_{\theta}(y\mid x)}\bigg[\underbrace{r_{\phi}(x, y) -\beta\log\sum_{y}\pi_\text{ref}(y\mid x)\exp\left(\frac{1}{\beta}r_{\phi}(x, y)\right)}_{f(r_{\phi}, \pi_\text{ref}, \beta)} - \underbrace{\beta\log\frac{\pi_{\theta}(y\mid x)}{\pi_\text{ref}(y\mid x)}}_{\text{KL}}\bigg]
\end{equation}

This objective is identical to the one optimized in previous studies 
\citep{ziegler2020finetuning, stiennon2022learning, bai2022training, ouyang2022training}, which use the DPO-equivalent reward corresponding to the reward class of $r_{\phi}$. Within this framework, the normalization component in $f(r_{\phi}, \pi_\text{ref}, \beta)$ can be viewed as the soft value function for the reference policy $\pi_\text{ref}$. Although this term does not influence the optimal solution, its absence may cause the policy gradient of the objective to exhibit high variance, resulting in unstable learning. One approach to handle the normalization term involves employing a learned value function, though this too can pose optimization challenges. Alternatively, earlier works have normalized rewards by using a human completion baseline, effectively a single-sample Monte-Carlo approximation of the normalization term. In contrast, the DPO reparameterization provides a reward function that does not rely on any baselines.

\section{Experiments}
Here, we empirically assess DPO's capacity to directly train policies from preferences. Initially, in a controlled text-generation environment, we investigate: how effectively does DPO balance maximizing reward while minimizing the KL-divergence relative to the reference policy, compared to standard preference learning methods such as PPO? Subsequently, we test DPO's effectiveness on larger models and more challenging RLHF tasks, including summarization and dialogue. Our findings indicate that with minimal hyperparameter tuning, DPO generally matches or surpasses strong baselines like RLHF with PPO, as well as outperforming the best of $N$ sampled trajectories based on a learned reward function. Prior to presenting these results, we outline the experimental setup; further details are provided in Appendix~\ref{app:exp_details}.

\textbf{Tasks.} Our study investigates three distinct open-ended text generation tasks. In all cases, the algorithms learn a policy from a preference dataset $\mathcal{D}=\bigl\{x^{(i)}, y_w^{(i)}, y_l^{(i)}\bigr\}_{i=1}^N$. For the \textbf{controlled sentiment generation} task, $x$ represents a movie review prefix taken from the IMDb dataset \cite{maas-EtAl:2011:ACL-HLT2011}, and the policy must generate a continuation $y$ exhibiting positive sentiment. To enable controlled evaluation, we \textit{create} preference pairs by generating outputs evaluated using a pre-trained sentiment classifier, ensuring that $p(\text{positive}\mid x,y_w)>p(\text{positive}\mid x,y_l)$. For supervised fine-tuning (SFT), GPT-2-large is fine-tuned to convergence on the training partition of the IMDB reviews (additional details in App~\ref{app:sentiment_details}). In the \textbf{summarization} task, $x$ corresponds to a Reddit forum post, and the policy must generate a summary $y$ capturing the key points. Consistent with prior work, we employ the Reddit TL;DR summarization dataset \citep{volske-etal-2017-tl} together with human preference data collected by \citeauthor{stiennon2022learning}. The SFT model is fine-tuned on human-written summaries of forum posts\footnote{\url{https://huggingface.co/CarperAI/openai_summarize_tldr_sft}} using the TRLX framework \citep{leandro_von_werra_2023_7790115} for RLHF. The human preference data was obtained by \citeauthor{stiennon2022learning} using samples from a different but similarly fine-tuned SFT model. Lastly, in the \textbf{single-turn dialogue} task, $x$ is a user query, which may range from astrophysics questions to requests for relationship advice. The policy must generate an engaging and helpful response $y$ to the query. We utilize the Anthropic Helpful and Harmless dialogue dataset \citep{bai2022training}, comprising 170k human-assistant dialogues. Each conversation concludes with two candidate responses produced by a large (though unspecified) language model and a label indicating which response the human preferred. Since no pre-trained SFT model exists in this setting, we fine-tune an off-the-shelf language model solely on the preferred completions to create the SFT model.

\textbf{Evaluation.} Our experiments employ two distinct evaluation strategies. To assess how well each algorithm optimizes the constrained reward maximization objective, in the controlled sentiment generation scenario, we measure each algorithm by the frontier of achieved reward and KL-divergence from the reference policy; this frontier is calculable since we have access to the true reward function (a sentiment classifier). Conversely, in practical applications where the true reward function is unknown, we assess algorithms based on their \textit{win rate} compared to a baseline policy, utilizing GPT-4 as a stand-in for human evaluation of summary quality and response helpfulness in the summarization and single-turn dialogue contexts, respectively. For summarization, the baseline is the reference summaries from the test set; for dialogue, the baseline is the preferred response from the test dataset. Although prior work indicates LMs may serve as superior automated evaluators compared to traditional metrics \citep{Chen2023ExploringTU}, we perform a human study to validate the use of GPT-4 for evaluation in Sec.~\ref{sec:human-judgments}. Our findings show that GPT-4’s judgments strongly correlate with human evaluations, with agreement between humans and GPT-4 generally matching or exceeding inter-human annotator agreement.

\textbf{Methods.} Alongside DPO, we assess several existing methods for training language models to conform to human preferences. At a basic level, we examine zero-shot prompting with \textbf{GPT-J} \citep{gpt-j} for summarization and 2-shot prompting with \textbf{Pythia-2.8B} \citep{biderman2023pythia} for dialogue. We also evaluate the \textbf{SFT} model and \textbf{Preferred-FT}, which is fine-tuned using supervised learning on the chosen completion $y_w$ derived from either the SFT model (in controlled sentiment and summarization) or a generic LM (in single-turn dialogue). Another pseudo-supervised technique is \textbf{Unlikelihood}~\citep{welleck2019neural}, which optimizes the policy to maximize the likelihood of $y_w$ while \textit{minimizing} the likelihood of $y_l$; we apply an optional weighting coefficient $\alpha\in[0,1]$ on the unlikelihood component. We further include \textbf{PPO} \citep{schulman2017proximal} trained with a reward function inferred from preference data, as well as \textbf{PPO-GT}, an oracle variant that learns using the true reward function available in the controlled sentiment setting. For sentiment experiments, we implement two versions of PPO-GT: an off-the-shelf version \cite{leandro_von_werra_2023_7790115} and a modified version that normalizes rewards and fine-tunes hyperparameters for enhanced performance (these modifications are also applied to standard PPO with learned rewards). Lastly, we evaluate the \textbf{Best of $N$} baseline, which samples $N$ responses from the SFT model (or Preferred-FT for dialogue) and selects the highest-scoring response according to a reward model trained on the preference dataset. This effective approach separates the reward model quality from PPO optimization but is computationally infeasible for moderate $N$ due to the necessity of generating $N$ completions per query at test time.

\subsection{How effectively can DPO optimize the RLHF objective?}

The KL-constrained reward maximization objective commonly used in RLHF methods aims to balance maximizing reward while preventing the policy from straying too far from the reference policy. Consequently, when comparing different algorithms, it is important to consider both the reward attained and the KL divergence; obtaining marginally higher reward at the cost of substantially larger KL divergence is not always beneficial. Figure~\ref{fig:frontier-tldr-main} illustrates the reward-KL frontier across various algorithms in the sentiment task. We perform multiple training runs for each algorithm, varying the policy conservativeness hyperparameter in each run (target KL $\in\{3,6,9,12\}$ for PPO, $\beta \in \{0.05,0.1,1,5\}$, $\alpha\in\{0.05,0.1,0.5,1\}$ for unlikelihood, and different random seeds for preferred-FT). This sweep totals 22 runs. After every 100 training steps until convergence, we evaluate each policy on a set of test prompts, calculating the average reward under the true reward function as well as the average sequence-level KL\footnote{That is, the sum of per-timestep KL divergences.} relative to the reference policy $\text{KL}\left(\pi\mid \mid \pi_\text{ref}\right)$. We observe that DPO yields the most efficient frontier by far, achieving the highest reward while maintaining low KL. This finding is notable for several reasons. First, although DPO and PPO optimize the same objective, DPO is significantly more efficient; its reward/KL tradeoff strictly outperforms PPO. Second, DPO attains a superior frontier even compared to PPO with access to ground truth rewards (PPO-GT).

\subsection{Is DPO scalable to actual preference datasets?}
\label{sec:dpo-real-datasets}
Next, we assess the fine-tuning effectiveness of DPO on tasks involving summarization and single-turn dialogue. For summarization, automatic metrics like ROUGE often show weak alignment with human judgments~\citep{stiennon2022learning}, and previous studies have demonstrated that fine-tuning language models with PPO based on human preferences yields better summary quality. We compare different methods by generating completions on the test portion of the TL;DR summarization dataset and calculating the average win rate against reference completions from the test set. Completions for all approaches are sampled across temperatures ranging from 0.0 to 1.0, with the corresponding win rates displayed in Figure~\ref{fig:frontier-tldr-main} (right). DPO, PPO, and Preferred-FT all fine-tune the same GPT-J SFT model\footnote{\url{https://huggingface.co/CarperAI/openai_summarize_tldr_sft}}. Our findings show that DPO achieves a win rate of about 61\% at a temperature of 0.0, surpassing PPO’s best performance of roughly 57\% at its optimal sampling temperature of 0.0. Additionally, DPO attains a higher peak win rate than the best-performing $N$ baseline. We note that we did not extensively optimize the $\beta$ hyperparameter for DPO, so these results could potentially undervalue DPO’s capabilities. Furthermore, DPO demonstrates greater robustness to variations in sampling temperature compared to PPO, whose performance can fall to the level of the base GPT-J model at elevated temperatures. Preferred-FT does not show significant improvements over the SFT model. In Section~\ref{sec:human-judgments}, we also present a direct comparison of DPO and PPO using human evaluations, where DPO outputs sampled at temperature 0.25 were favored 58\% of the time over PPO outputs sampled at temperature 0.

For single-turn dialogue, we assess various approaches on a subset of the Anthropic HH dataset test split \citep{bai2022training} containing one round of human-assistant interaction. GPT-4 evaluations utilize the preferred completions in the test set as a reference to calculate the win rate across different methods. Since no standard SFT model exists for this task, we begin with a pre-trained Pythia-2.8B, apply Preferred-FT to train a reference model on the selected completions to ensure the completions align with the model's distribution, and subsequently train using DPO. We also benchmark against the top completion among 128 Preferred-FT outputs (noting that the Best of $N$ baseline plateaus at 128 completions for this task; see Appendix Figure~\ref{fig:best-of-n}) and a 2-shot prompted version of the Pythia-2.8B base model, observing that DPO matches or exceeds the best-performing temperature results for each method. Additionally, we evaluate an RLHF model trained via PPO on the Anthropic HH dataset \footnote{\url{https://huggingface.co/reciprocate/ppo_hh_pythia-6B}} from a well-known source \footnote{\url{https://github.com/CarperAI/trlx/tree/main/examples/hh}}, but cannot identify a prompt or sampling temperature that outperforms the base Pythia-2.8B model. Drawing from our TL;DR results and noting that both methods optimize the identical reward function, we treat Best of 128 as an approximate proxy for PPO-level performance. In summary, DPO stands out as the only computationally efficient approach that surpasses the preferred completions on the Anthropic HH dataset, achieving comparable or better outcomes than the computationally intensive Best of 128 baseline. Lastly, Figure~\ref{fig:dialogue-main} illustrates that DPO attains its optimal performance relatively rapidly.

\subsection{Generalization to a novel input distribution}

To further assess how PPO and DPO perform under shifts in data distribution, we test the PPO and DPO policies from our Reddit TL;DR summarization experiment on a different dataset: news articles from the CNN/DailyMail test split \citep{nallapati-etal-2016-abstractive}, using the optimal sampling temperatures derived from TL;DR (0 and 0.25). The outcomes are shown in Table~\ref{tab:ood}. We calculated the GPT-4 win rate compared to the ground-truth summaries in these datasets, employing the same GPT-4 (C) prompt used for Reddit TL;DR, but substituting the term ``forum post'' with ``news article''. For this new distribution, DPO continues to significantly outperform the PPO policy. This experiment offers preliminary evidence that DPO policies can generalize just as effectively as PPO policies, despite DPO not leveraging the additional unlabeled Reddit TL;DR prompts utilized by PPO.

\subsection{Validating GPT-4 evaluations against human judgments}
\label{sec:human-judgments}
We conducted a human evaluation study to confirm the reliability of GPT-4’s assessments, using results from the TL;DR summarization experiment and two distinct GPT-4 prompts. The \textbf{GPT-4 (S)} (simple) prompt straightforwardly asks which summary better captures the key information in the post. The \textbf{GPT-4 (C)} (concise) prompt additionally requests which summary is more concise; we test this prompt because we observe that GPT-4 tends to favor longer, more repetitive summaries than humans do when using the \textbf{GPT-4 (S)} prompt. Complete prompt details are in Appendix~\ref{app:prompts}. We performed three pairwise comparisons involving the highest-performing method (DPO, temp. 0.25), the lowest-performing method (PPO, temp. 1.0), and a middle-performing method (SFT, temp. 0.25), with the goal of capturing a range of sample qualities; all three were compared against greedily sampled PPO (its best temperature). Our findings show that GPT-4 agrees with human raters about as frequently as humans agree with one another, indicating that GPT-4 serves as a reasonable proxy for human evaluation (due to a limited number of human raters, multiple human judgments were collected only for the DPO and PPO-1 comparisons). Overall, the \textbf{GPT-4 (C)} prompt tends to yield win rates more closely aligned with human preferences; therefore, we adopt this prompt for the main results reported in Section~\ref{sec:dpo-real-datasets}. Additional information about the human study, including the web interface shown to raters and a list of human participants, is available in Appendix~\ref{app:human-study}.

\section{Discussion}
Learning from preferences offers a robust and scalable approach for developing capable and aligned language models. We have presented DPO, a straightforward training methodology that enables language models to be trained from preferences without relying on reinforcement learning. Instead of forcing the preference learning problem into a conventional RL framework to utilize standard RL algorithms, DPO establishes a direct correspondence between language model policies and reward functions. This allows training a language model to align with human preferences \textit{directly}, using a simple cross-entropy loss, without the need for reinforcement learning or sacrificing generality. With minimal hyperparameter tuning, DPO achieves performance comparable to or surpassing existing RLHF methods, including those based on PPO; thus, DPO significantly lowers the barrier to training additional language models guided by human preferences.

\textbf{Limitations \& Future Work.} Our findings bring up several important avenues for further investigation. How well does the DPO policy generalize to out-of-distribution data compared to learning from an explicit reward function? Preliminary results indicate that DPO policies generalize on par with PPO-based models, but a more thorough analysis is warranted. For instance, can self-labeling approaches with the DPO policy effectively utilize unlabeled prompts? Moreover, how does reward over-optimization present itself within the direct preference optimization framework, and could the slight performance drop shown in Figure~\ref{fig:dialogue-main}-right be an example of this? While we evaluate models up to 6B parameters, expanding DPO to scale to state-of-the-art models that are orders of magnitude larger represents a promising direction for future research. Regarding evaluations, we observe that GPT-4’s win rate assessments are influenced by the prompt, suggesting that future work might explore optimal methods for eliciting high-quality judgments from automated evaluators. Lastly, the potential applications of DPO extend well beyond training language models from human preferences, including training generative models in various other modalities.

\end{document}